\chapter{机械振子与时间的测量}\label{ch:g12:oscillators-and-time}

落地钟滴答作响：每一次滴答，都是黄铜摆锤经过
竖直方向，自元旦以来已重复一千五百万次。
你手表里的石英薄片每秒振动 \num{32768} 次；铯\omterm{def:g11:fundamental-interactions:ladder}{原子}
振动得快九十亿倍。凡能摆动的东西都能计时：
本章循着摆动的线索，从惠更斯讲到\omterm{def:g10:universal-gravitation:satellite}{卫星}。

\section{自由振动}

\begin{definition}[振子、周期、频率、振幅]\label{def:g12:oscillators-and-time:oscillator}
\emph{振子}\index{振子}是这样一种系统：
它被从稳定的\emph{平衡位置}\index{平衡位置}处移开后释放，
会绕该位置来回摆动，这种运动称为
\emph{自由振动}\index{自由振动}。运动每隔一个
\emph{周期}\index{周期} $T_0$（\unit{s}）重复一次；
\emph{频率}\index{频率} $f_0 = 1/T_0$
（\unit{Hz}，\cref{ch:g12:mechanical-waves}）；
\emph{振幅}\index{振幅}是最大位移。
\end{definition}

\begin{example}[形形色色的振子]\label{ex:g12:oscillators-and-time:zoo}
游乐场的秋千：$T_0 \approx \qty{3}{s}$。吉他的 A 弦：
$f_0 = \qty{110}{Hz}$。石英晶体：\qty{32768}{Hz}。定义秒的铯\omterm{def:g11:fundamental-interactions:ladder}{原子}
振动：约 \qty{9.2e9}{Hz}。跨越十个
\omterm{def:g10:orders-of-magnitude:oom}{数量级}，却只做一件事：计数。
\end{example}

\section{单摆}

\begin{definition}[单摆]\label{def:g12:oscillators-and-time:pendulum}
\emph{单摆}\index{单摆}：一个小摆球，质量为
$m$，系在长为 $L$、质量可忽略的不可伸长细线上，
在竖直平面内摆动。摆球受两个力
（\cref{ch:g12:newtons-laws}）：\omterm{def:g10:inertia:inventory}{张力}与运动方向
垂直，而\omterm{def:g10:universal-gravitation:weight}{重力}沿\omterm{def:g10:relative-motion:trajectory}{轨迹}方向的分量则指向最低点，
因此\omterm{def:g10:universal-gravitation:weight}{重力}就是\emph{回复力}%
\index{回复力}，与位移方向相反。
\end{definition}

\begin{omfigure}
\begin{tikzpicture}[scale=0.8, font=\small]
  \coordinate (O) at (0,0); \draw (-0.6,0) -- (0.6,0);
  \foreach \x in {-0.5,-0.25,0,0.25,0.5} \draw (\x,0) -- ++(0.14,0.14);
  \fill (O) circle (1.5pt); \draw[dashed] (O) -- (0,-3.4);
  \draw[dashed, omExo] (-120:3.0) arc (-120:-60:3.0);
  \draw (O) -- (-60:3.0) node[midway, above right] {$L$};
  \draw (-90:1.0) arc (-90:-60:1.0); \node at (-75:1.32) {$\theta$};
  \coordinate (B) at (1.5,-2.598); \fill[omDef] (B) circle (2.5pt);
  \draw[-Stealth, omDef, thick] (B) -- ++(-0.75,1.299) node[above left] {$\vect T$};
  \draw[-Stealth, thick] (B) -- ++(0,-1.8) node[below] {$\vect P$};
  \draw[-Stealth, omThm, thick] (B) -- ++(-0.779,-0.45) node[left] {$\vect{P_t}$};
  \draw[dashed, omExo] (0.721,-3.048) -- (1.5,-4.398)
    (B) -- ++(0.779,-1.35) (2.279,-3.948) -- (1.5,-4.398);
\end{tikzpicture}

{\small 摆球受力：\omterm{def:g10:inertia:inventory}{张力} $\vect T$ 与
运动方向垂直；\omterm{def:g10:universal-gravitation:weight}{重力}的切向分量 $\vect{P_t}$ 总是
指向竖直方向——即\omterm{def:g12:oscillators-and-time:pendulum}{回复力}。}
\end{omfigure}

\begin{proposition}[单摆的周期]\label{prop:g12:oscillators-and-time:pendulum-period}
\index{单摆!周期}
小\omterm{def:g12:oscillators-and-time:oscillator}{振幅}（约 $15^\circ$ 以下）时，\omterm{def:g12:oscillators-and-time:oscillator}{周期}为
\[
T_0 = 2\pi \sqrt{\frac{L}{g}}, \qquad g = \qty{9.81}{m/s^2}.
\]
\omterm{def:g12:oscillators-and-time:oscillator}{周期}与摆球的质量无关，并且在角度保持较小时，
也与\omterm{def:g12:oscillators-and-time:oscillator}{振幅}无关：小幅摆动具有
\emph{等时性}\index{等时性}，各次摆动所用时间相同。
\end{proposition}
\admitted

\begin{remark}[量纲分析，以及 $2\pi$ 从何而来]\label{rem:g12:oscillators-and-time:dimensional}
这个公式可以猜出。由 $m$（\unit{kg}）、$L$（\unit{m}）和
$g$（\unit{m/s^2}）出发，唯一具有时间\omterm{def:g10:orders-of-magnitude:unit}{单位}的组合是
$\sqrt{L/g}$——而质量\emph{不能}进入，
因为没有别的量携带\omterm{def:g10:orders-of-magnitude:unit}{千克}\omterm{def:g10:orders-of-magnitude:unit}{单位}能与之抵消。量纲分析无法给出
前面的纯数：$2\pi$ 来自对运动方程的求解，
这一求解将在本科一年级分册中完成。小角度条件也不可省略：
在 $30^\circ$ 时，真实\omterm{def:g12:oscillators-and-time:oscillator}{周期}要长 $2\%$，\omterm{prop:g12:oscillators-and-time:pendulum-period}{等时性}就不再成立。
\end{remark}

\begin{example}[秒摆]\label{ex:g12:oscillators-and-time:seconds}
哪一个摆长能“打出秒”，使\omterm{def:g12:oscillators-and-time:oscillator}{周期} $T_0 = \qty{2.00}{s}$，
即每个方向各滴答一次？求解得 $L = g\,T_0^2/4\pi^2 = 9.81 \times 2.00^2/4\pi^2 \approx
\qty{0.994}{m}$。接近一米：摆钟之所以高大，
是因为秒本身较长。
\end{example}

\section{弹簧–振子}

\begin{definition}[弹簧的回复力]\label{def:g12:oscillators-and-time:spring}
质量为 $m$ 的滑块在水平无摩擦导轨上滑动，并与
一根\emph{劲度系数}\index{劲度系数}为 $k$（\unit{N/m}）的弹簧相连。设
$x$ 为其离开\omterm{def:g12:oscillators-and-time:oscillator}{平衡位置}的位移。被拉伸或压缩时，
弹簧以力 $F = -kx$ 将滑块拉回或推回：
力与位移成正比、方向与位移相反——又一个
\omterm{def:g12:oscillators-and-time:pendulum}{回复力}，这次由弹性提供。
\end{definition}

\begin{proposition}[运动方程及其解]\label{prop:g12:oscillators-and-time:spring-period}
\index{弹簧–振子}
沿导轨方向的\omterm{def:g10:inertia:force}{牛顿}第二定律（\cref{ch:g12:newtons-laws}）给出
$m \frac{d^2x}{dt^2} = -kx$，即 $\frac{d^2x}{dt^2} =
-\frac{k}{m}\,x$。无论常数 $A$ 和 $\varphi$ 取何值，
\[
x(t) = A \cos\!\left(\frac{2\pi t}{T_0} + \varphi\right),
\qquad T_0 = 2\pi\sqrt{\frac{m}{k}},
\]
满足该方程：\omterm{def:g12:oscillators-and-time:oscillator}{周期}为 $T_0$ 的振动，与
\omterm{def:g12:oscillators-and-time:oscillator}{振幅}无关。
\end{proposition}

\begin{proof}
求导，得 $\frac{dx}{dt} = -\frac{2\pi}{T_0}A \sin(\frac{2\pi
t}{T_0} + \varphi)$；再求一次导，$\frac{d^2x}{dt^2} = -(\frac{2\pi}{T_0})^2
A\cos(\frac{2\pi t}{T_0} + \varphi) = -(\frac{2\pi}{T_0})^2\,x(t)$,
即 $-(k/m)\,x$，当且仅当 $(2\pi/T_0)^2 = k/m$ 时成立。
\end{proof}

\begin{remark}[所有可能的运动]\label{rem:g12:oscillators-and-time:uniqueness}
滑块的\emph{每一种}运动都是这些函数之一，这一点
在本阶段不加证明地承认，其证明见本科一年级分册。另外，$g$ 不出现在 $2\pi\sqrt{m/k}$ 中：
\omterm{def:g12:oscillators-and-time:spring}{劲度系数}越大，振动越快；质量越大，振动越慢；在月球上也是如此。
\end{remark}

\begin{definition}[振幅与相位]\label{def:g12:oscillators-and-time:phase}
在 $x(t) = A\cos(2\pi t/T_0 + \varphi)$ 中，正常数 $A$ 称为
\omterm{def:g12:oscillators-and-time:oscillator}{振幅}——最大位移处为 $x = \pm A$——而 $2\pi t/T_0 +
\varphi$ 称为\emph{相位}\index{相位}，每过一个\omterm{def:g12:oscillators-and-time:oscillator}{周期}
增加 $2\pi$；\emph{初相位}\index{初相位} $\varphi$
记录起始状态：自 $x = A$ 处静止释放时，$\varphi = 0$。例如，
当 $m = \qty{0.20}{kg}$、$k = \qty{20}{N/m}$ 时，得 $T_0 =
2\pi\sqrt{0.20/20} = \qty{0.63}{s}$；将其拉至 \qty{5.0}{cm} 处，
释放后 $x(t) = 5.0\cos(2\pi t/0.63)$，\omterm{def:g10:orders-of-magnitude:unit}{单位}为厘米。
\end{definition}

\section{能量、阻尼与共振}

\begin{proposition}[振子的能量]\label{prop:g12:oscillators-and-time:energy}
伸长量为 $x$ 的弹簧储存的弹性势能为 $E_p =
\tfrac12 kx^2$（依据见\cref{ch:g12:work-and-energy}）。无摩擦
的振动中，\omterm{def:g11:mechanical-energy:em}{机械能}
\[
E_m = E_k + E_p = \tfrac12 m v^2 + \tfrac12 k x^2 = \tfrac12 k A^2
\]
保持不变：最大位移处全部为弹性势能，\omterm{def:g12:oscillators-and-time:oscillator}{平衡位置}处全部为\omterm{def:g11:mechanical-energy:kinetic}{动能}。
\omterm{def:g12:oscillators-and-time:pendulum}{单摆}在 $E_k$ 与\omterm{def:g11:mechanical-energy:potential}{重力势能}
$E_p$ 之间进行同样的\omterm{def:g10:energy-conservation:energy}{能量}转换（见\cref{ch:g11:mechanical-energy}）。
\end{proposition}

\begin{proof}
记 $v = \frac{dx}{dt}$，则 $E_m$ 的导数为 $mv\frac{dv}{dt}
+ kx\frac{dx}{dt} = v\,(m\frac{d^2x}{dt^2} + kx) = 0$，由运动
方程；数值 $\tfrac12 kA^2$ 在最大位移处读出。
\end{proof}

\begin{omfigure}
\begin{tikzpicture}
\begin{axis}[omaxis, width=8.6cm, height=4.6cm, xlabel={$t/T_0$},
    ylabel={能量 $/\,E_m$}, xmin=0, xmax=1, ymin=0, ymax=1.3,
    legend pos=south east]
  \addplot[omDef, thick, domain=0:1, samples=120] {cos(360*x)^2};
  \addplot[omProp, thick, domain=0:1, samples=120] {sin(360*x)^2};
  \addplot[omThm, dashed, thick, domain=0:1, samples=2] {1};
  \legend{$E_p$, $E_k$, $E_m$}
\end{axis}
\end{tikzpicture}

{\small 一个\omterm{def:g12:oscillators-and-time:oscillator}{周期}内，从 $x = A$ 处释放：\omterm{def:g10:energy-conservation:energy}{能量}在弹性势能与\omterm{def:g11:mechanical-energy:kinetic}{动能}间往返两次，
总量恒为 $\tfrac12 kA^2$。}
\end{omfigure}

\begin{definition}[阻尼振动]\label{def:g12:oscillators-and-time:damping}
\omterm{def:g10:inertia:inventory}{摩擦力}会从每个实际\omterm{def:g12:oscillators-and-time:oscillator}{振子}中\omterm{def:g11:mechanical-energy:dissipation}{耗散}\omterm{def:g11:mechanical-energy:em}{机械能}：其
\omterm{def:g12:oscillators-and-time:oscillator}{自由振动}为\emph{阻尼}\index{阻尼}振动。在轻阻尼
情形下，系统仍做振动，\omterm{def:g12:oscillators-and-time:oscillator}{振幅}逐渐\omterm{def:g12:mechanical-waves:attenuation}{衰减}，但
\omterm{def:g12:oscillators-and-time:oscillator}{周期}几乎不变，此即\emph{准\omterm{def:g12:oscillators-and-time:oscillator}{周期}}%
\index{准周期}，接近 $T_0$；在重阻尼（如蜂蜜中）
系统缓慢回到\omterm{def:g12:oscillators-and-time:oscillator}{平衡位置}，而不再越过它；\omterm{def:g10:energy-conservation:energy}{能量}
与\omterm{def:g12:oscillators-and-time:oscillator}{振幅}的平方成正比，最终\omterm{def:g11:mechanical-energy:dissipation}{耗散}为热量。
\end{definition}

\begin{omfigure}
\begin{tikzpicture}
\begin{axis}[omaxis, width=8.6cm, height=4.6cm, xlabel={$t$ (\unit{s})},
    ylabel={$x$ (\unit{cm})}, xmin=0, xmax=10, ymin=-2.6, ymax=2.6,
    legend pos=south east]
  \addplot[omExo, thin, domain=0:10, samples=240] {2*cos(deg(pi*x))};
  \addplot[omDef, thick, domain=0:10, samples=240] {2*exp(-x/4)*cos(deg(pi*x))};
  \legend{无阻尼, 有阻尼}
  \addplot[omThm, dashed, domain=0:10, samples=60, forget plot] {2*exp(-x/4)};
  \addplot[omThm, dashed, domain=0:10, samples=60, forget plot] {-2*exp(-x/4)};
\end{axis}
\end{tikzpicture}

{\small 无\omterm{def:g12:oscillators-and-time:damping}{阻尼}与轻\omterm{def:g12:oscillators-and-time:damping}{阻尼}振动：\omterm{def:g12:oscillators-and-time:oscillator}{周期}几乎相同，
（此处为 \qty{2.0}{s}），\omterm{def:g12:oscillators-and-time:damping}{阻尼}振动的\omterm{def:g12:oscillators-and-time:oscillator}{振幅}
在逐渐收缩的\omterm{def:g10:orders-of-magnitude:scinot}{指数}包络线（虚线）内\omterm{def:g12:mechanical-waves:attenuation}{衰减}。}
\end{omfigure}

\begin{definition}[受迫振动与共振]\label{def:g12:oscillators-and-time:resonance}
对\omterm{def:g12:oscillators-and-time:oscillator}{振子}施加\omterm{def:g12:oscillators-and-time:oscillator}{周期}性推动——即\emph{受迫振动}%
\index{受迫振动}——经过短暂过渡后，它
以\emph{驱动}\omterm{def:g12:oscillators-and-time:oscillator}{频率} $f$ 振动，而非其自身的\omterm{def:g12:oscillators-and-time:oscillator}{频率}。
其\omterm{def:g12:oscillators-and-time:oscillator}{振幅}在 $f$ 接近自由\omterm{def:g12:oscillators-and-time:oscillator}{振子}的\emph{固有\omterm{def:g12:oscillators-and-time:oscillator}{频率}}%
\index{固有频率} $f_0$ 时达到峰值：此即
\emph{共振}\index{共振}：每次推动都与
运动同步，每个\omterm{def:g12:oscillators-and-time:oscillator}{周期}都输入\omterm{def:g10:energy-conservation:energy}{能量}；\omterm{def:g12:oscillators-and-time:damping}{阻尼}越轻，
峰值越高、越尖锐。
\end{definition}

\begin{omfigure}
\begin{tikzpicture}
\begin{axis}[omaxis, width=8.6cm, height=4.6cm, xlabel={$f/f_0$},
    ylabel={振幅（任意单位）}, xmin=0, xmax=2.2, ymin=0, ymax=5.6,
    legend pos=north east]
  \addplot[omThm, thick, domain=0:2.2, samples=200] {1/sqrt((1-x^2)^2 + (x/5)^2)};
  \addplot[omDef, thick, domain=0:2.2, samples=120] {1/sqrt((1-x^2)^2 + (x/1.2)^2)};
  \legend{轻阻尼, 重阻尼}
  \addplot[omExo, dashed, forget plot] coordinates {(1,0) (1,5.4)};
\end{axis}
\end{tikzpicture}

{\small \omterm{def:g12:oscillators-and-time:resonance}{共振}曲线：\omterm{def:g12:oscillators-and-time:resonance}{受迫振动}的\omterm{def:g12:oscillators-and-time:oscillator}{振幅}随驱动\omterm{def:g12:oscillators-and-time:oscillator}{频率}变化的关系。
在 $f = f_0$ 附近响应急剧增大，\omterm{def:g12:oscillators-and-time:damping}{阻尼}越轻，峰值越高。}
\end{omfigure}

\begin{example}[有利与有害的共振]\label{ex:g12:oscillators-and-time:resonance-zoo}
家长推秋千时，总以秋千自身的\omterm{def:g12:oscillators-and-time:oscillator}{频率}用力——这正是让\omterm{def:g12:oscillators-and-time:resonance}{共振}为人所用。
\num{2000} 年，伦敦新建的千禧桥摇晃得令人担忧，
原因是人群的步伐锁定在接近
\qty{1}{Hz} 的\omterm{def:g12:oscillators-and-time:resonance}{固有频率}上，桥因此关闭两年以加装\omterm{def:g12:oscillators-and-time:damping}{阻尼}器。每只石英表都靠
以\omterm{def:g12:oscillators-and-time:resonance}{共振}方式驱动石英晶体持续振动，\omterm{def:g12:oscillators-and-time:oscillator}{频率}为 \qty{32768}{Hz}。
\end{example}

\section{测量时间}

\begin{method}[周期的测定]\label{met:g12:oscillators-and-time:timing}
不要只测一次振动：在摆球经过竖直方向时开始按秒表
（此时最快、最易判断），数出 $N = 50$ 次振动，
再除以 $N$——一个 \qty{0.2}{s} 的反应误差，
在 $T_0$ 上变为 \qty{4}{ms}。重复测量，并比较各次结果。
\end{method}

\begin{example}[三个世纪的滴答]\label{ex:g12:oscillators-and-time:clocks}
惠更斯于 \num{1657} 年制成第一台摆钟：\omterm{prop:g12:oscillators-and-time:pendulum-period}{等时性}
使钟表的日误差由每天一刻钟降到秒级，而
补偿摆则把日误差降到不足一秒。石英
腕表（\num{1969} 年）以石英晶体取代摆杆，晶体以
\qty{32768}{Hz} 振动：月误差仅十分之几秒。自 \num{1967} 年起，
秒本身改由一次\omterm{def:g11:fundamental-interactions:ladder}{原子}振荡来\emph{定义}：秒即
铯-133 \omterm{def:g11:fundamental-interactions:ladder}{原子}内部跃迁所辐射微波的 \num{9192631770} 个\omterm{def:g12:oscillators-and-time:oscillator}{周期}的时长
（见\cref{ch:g12:quantum-world}）。最好的铯\omterm{def:g11:fundamental-interactions:ladder}{原子}钟
一亿年才偏差一秒。
\end{example}

\begin{remark}[为何追求纳秒精度]\label{rem:g12:oscillators-and-time:gps}
GPS 接收机通过测量无线电\omterm{def:g10:signals-and-waves:signal}{信号}的时间来确定位置，这些\omterm{def:g10:signals-and-waves:signal}{信号}
来自\omterm{def:g10:universal-gravitation:satellite}{卫星}；光每纳秒传播 \qty{30}{cm}，因此
时钟误差一微秒，定位就会偏差 \qty{300}{m}。\omterm{def:g10:universal-gravitation:satellite}{卫星}
因此携带\omterm{def:g11:fundamental-interactions:ladder}{原子}钟——在这一精度下，时钟的快慢
取决于它运动的快慢与所处的高度，这种效应
将在\cref{ch:g12:special-relativity}中讨论。
\end{remark}

\section{练习}

\begin{exercise}[$\star$]\label{exo:g12:oscillators-and-time:1}
(a) 求长为 \qty{1.00}{m} 的\omterm{def:g12:oscillators-and-time:pendulum}{单摆}的\omterm{def:g12:oscillators-and-time:oscillator}{周期}与\omterm{def:g12:oscillators-and-time:oscillator}{频率}。
(b) 摆长变为原来的四倍，\omterm{def:g12:oscillators-and-time:oscillator}{周期}将如何变化？
\end{exercise}

\begin{exercise}[$\star$]\label{exo:g12:oscillators-and-time:2}
静止时心脏每分钟跳动 $72$ 次；蜻蜓翅膀以
\qty{30}{Hz} 的\omterm{def:g12:oscillators-and-time:oscillator}{频率}振动。求心脏的\omterm{def:g12:oscillators-and-time:oscillator}{频率}与\omterm{def:g12:oscillators-and-time:oscillator}{周期}，以及翅膀的
\omterm{def:g12:oscillators-and-time:oscillator}{周期}；本章中哪个词表示每次搏动时胸部的最大起伏？
\end{exercise}

\begin{exercise}[$\star$]\label{exo:g12:oscillators-and-time:3}
一个质量为 \qty{0.10}{kg} 的滑块，装在\omterm{def:g12:oscillators-and-time:spring}{劲度系数}为 \qty{25}{N/m} 的弹簧上：求 $T_0$
和 $f_0$。质量加倍，\omterm{def:g12:oscillators-and-time:oscillator}{周期}将如何变化？
\end{exercise}

\begin{exercise}[$\star$]\label{exo:g12:oscillators-and-time:4}
\omterm{def:g12:oscillators-and-time:damping}{阻尼}\omterm{def:g12:oscillators-and-time:oscillator}{振子}相继的极大值分别为 $2.0$、$1.5$、$1.13$，
\qty{0.85}{cm}，对应时刻 $t = 0$、$2.0$、$4.0$、\qty{6.0}{s}。读出
\omterm{def:g12:oscillators-and-time:damping}{准周期}；说明每经过一个\omterm{def:g12:oscillators-and-time:oscillator}{周期}，\omterm{def:g12:oscillators-and-time:oscillator}{振幅}都按同一固定比例减小；
并预测下一个极大值。
\end{exercise}

\begin{exercise}[$\star$]\label{exo:g12:oscillators-and-time:5}
一名宇航员把一个\omterm{def:g12:oscillators-and-time:pendulum}{单摆}和一个弹簧–质量\omterm{def:g12:oscillators-and-time:oscillator}{振子}带到
月球（$g = \qty{1.62}{m/s^2}$）上，两种\omterm{def:g12:oscillators-and-time:oscillator}{周期}各变为原来的多少倍？
并说明理由。
\end{exercise}

\begin{exercise}[$\star\star$]\label{exo:g12:oscillators-and-time:6}
(a) 验证 $\sqrt{L/g}$ 具有时间的量纲。(b)
为何 $T_0$ 的公式不可能含有质量，哪怕它只由 $m$、$L$、$g$ 构成？
(c) 这一论证能否得出 $2\pi$？
\end{exercise}

\begin{exercise}[$\star\star$]\label{exo:g12:oscillators-and-time:7}
通过两次求导验证 $x(t) = A\cos(2\pi t/T_0 +
\varphi)$ 满足 $\frac{d^2x}{dt^2} = -(k/m)\,x$ 当且仅当
$T_0 = 2\pi\sqrt{m/k}$。
\end{exercise}

\begin{exercise}[$\star\star$]\label{exo:g12:oscillators-and-time:8}
一根\omterm{def:g12:oscillators-and-time:spring}{劲度系数}为 \qty{40}{N/m} 的弹簧，连着一个 \qty{0.20}{kg} 的滑块，
\omterm{def:g12:oscillators-and-time:oscillator}{振幅}为 \qty{3.0}{cm}。求\omterm{def:g11:mechanical-energy:em}{机械能}与
最大速率，并指出速率在何处最大、何处为零？
\end{exercise}

\begin{exercise}[$\star\star$]\label{exo:g12:oscillators-and-time:9}
(a) 求秒摆的摆长（$T_0 =
\qty{2.00}{s}$）。(b) 若摆长增加 $1.0\%$，利用 $\sqrt{1+x}
\approx 1 + x/2$，求 $T_0$ 的相对变化及日误差。
\end{exercise}

\begin{exercise}[$\star\star$]\label{exo:g12:oscillators-and-time:10}
一个轻\omterm{def:g12:oscillators-and-time:damping}{阻尼}\omterm{def:g12:oscillators-and-time:oscillator}{振子}每个\omterm{def:g12:oscillators-and-time:oscillator}{周期}损失其\omterm{def:g10:energy-conservation:energy}{能量}的 $5.0\%$。
(a) 每个\omterm{def:g12:oscillators-and-time:oscillator}{周期}\omterm{def:g12:oscillators-and-time:oscillator}{振幅}下降多少？(b) 损失一半的
\omterm{def:g10:energy-conservation:energy}{能量}需要多少个\omterm{def:g12:oscillators-and-time:oscillator}{周期}？(c) \omterm{def:g12:oscillators-and-time:damping}{准周期}是否与 $T_0$ 相差甚远？
\end{exercise}

\begin{exercise}[$\star\star$]\label{exo:g12:oscillators-and-time:11}
借助\omterm{def:g12:oscillators-and-time:resonance}{共振}曲线说明：(a) 为何推秋千只在
其自身节奏时才有效；(b) 为何士兵过人行桥时要便步；(c)
为何洗衣机在升速甩干时会在某一转速下剧烈振动。
\end{exercise}

\begin{exercise}[$\star\star\star$]\label{exo:g12:oscillators-and-time:12}
一台在海平面走时精确的摆钟被移到一座天文台，那里
$g = \qty{9.78}{m/s^2}$。(a) 走快还是走慢？每
天相差多少秒？(b) 它那根 \qty{0.994}{m} 的摆应缩短多少？
\end{exercise}

\begin{exercise}[$\star\star\star$]\label{exo:g12:oscillators-and-time:13}
一个质量为 \qty{0.20}{kg} 的滑块满足 $x(t) = A\cos(2\pi t/T_0)$，其中 $A =
\qty{4.0}{cm}$，$T_0 = \qty{0.63}{s}$。求最大速率、
最大加速度、\omterm{def:g12:oscillators-and-time:spring}{劲度系数} $k$ 以及\omterm{def:g11:mechanical-energy:em}{机械能}。
\end{exercise}

\begin{exercise}[$\star\star\star$]\label{exo:g12:oscillators-and-time:14}
为测量 $g$：一个摆长 $L = 99.5 \pm \qty{0.2}{cm}$ 的\omterm{def:g12:oscillators-and-time:pendulum}{单摆}，
完成 $20$ 次振动用时 $40.1 \pm \qty{0.2}{s}$。求 $g$ 及其
\omterm{def:g10:orders-of-magnitude:uncertainty}{相对不确定度}（将 $L$ 与 $T_0^2$ 的\omterm{def:g10:orders-of-magnitude:uncertainty}{相对不确定度}相加）；
结果是否与 \qty{9.81}{m/s^2} 相符？
\end{exercise}

\begin{exercise}[$\star\star\star$]\label{exo:g12:oscillators-and-time:15}
用相对漂移（误差与经过时间之比）比较计时器：一台摆钟每天慢
\qty{10}{s}，一块石英表每月慢 \qty{0.5}{s}，
一台铯钟每一亿年慢 \qty{1}{s}，
求这三个比值。一台误差为
\qty{1}{\micro s} 的 GPS 时钟，其位置误差为多少？GPS 需要哪一类时钟？
\end{exercise}

\section{周末问题：钟表匠的工作台}

\begin{problem}[{周末问题——钟表匠的工作台：
一只挂钟被调回准时——测定五十次摆动，黄铜杆夏季伸长，
擒纵机构的轻推，一块石英挑战者——
结论以每月秒数表示}]
\label{pb:g12:oscillators-and-time:1}
一台挂钟来到工作台时“走时不准”。把它的黄铜杆和沉重的摆锤
看作等效摆长为 $L$ 的\omterm{def:g12:oscillators-and-time:pendulum}{单摆}，摆锤质量 $m = \qty{1.2}{kg}$；
齿轮系统把每次完整振动计为
恰好 \qty{2.000}{s}。

\medskip
\noindent\textbf{第一部分——测量挂钟的“脉搏”。}
\begin{enumerate}
  \item 为何要测 $50$ 次振动而非一次？若你的反应时间为 \qty{0.2}{s}，
        估计由此在 $T_0$ 上留下的误差。
  \item 秒表显示 $50$ 次振动用时 \qty{99.4}{s}，其\omterm{def:g12:oscillators-and-time:oscillator}{周期}为多少？
  \item 推算摆的当前等效摆长。
  \item 摆长为多少时，\omterm{def:g12:oscillators-and-time:oscillator}{周期}恰好为 \qty{2.000}{s}？
        \omterm{def:g11:lenses-and-eye:eye}{调节}螺母应使摆锤上移还是下移，移动多少？
  \item 不加\omterm{def:g11:lenses-and-eye:eye}{调节}，挂钟每天是快还是慢？相差多少秒？
  \item 对 $g$ 的灵敏度：月球上的\omterm{def:g12:oscillators-and-time:oscillator}{周期}为多少
        （$g = \qty{1.62}{m/s^2}$）？在 $g = \qty{9.80}{m/s^2}$ 的天文台，
        每天误差为多少？
\end{enumerate}

\medskip
\noindent\textbf{第二部分——夏季的“走慢”。}黄铜会膨胀：温度升高 $\Delta\theta$ 时，
黄铜杆伸长为 $L' = L\,(1 + \lambda\,
\Delta\theta)$，$\lambda = \num{1.9e-5}$（每\omterm{def:g10:light-spectra:kelvin}{开尔文}）。挂钟在
\qty{20}{\celsius} 时调准。
\begin{enumerate}[resume]
  \item 求夏季室温为 \qty{25}{\celsius} 时的 $\Delta L$。
  \item 利用 $\sqrt{1+x} \approx 1 + x/2$，证明
        $\Delta T_0 / T_0 = \tfrac12 \lambda\, \Delta\theta$。
  \item 求 $\Delta T_0/T_0$ 及夏季每日误差，并判断是快还是慢？
  \item 对冬季走廊（\qty{10}{\celsius}）的情况，回答同样的问题。
  \item 旧式精密挂钟采用“格架摆”，用黄铜杆与钢杆
        交替排列组成（钢杆膨胀约为黄铜杆的一半）。
        试说明这种结构如何保持走时稳定。
\end{enumerate}

\medskip
\noindent\textbf{第三部分——耳语般的推力。}擒纵机构脱开时，
以\omterm{def:g12:oscillators-and-time:oscillator}{振幅} $\theta_m = 2.00^\circ$ 开始摆动的摆，其\omterm{def:g12:oscillators-and-time:oscillator}{振幅}在 \qty{280}{s} 内减半；
正常工作时，擒纵机构每个\omterm{def:g12:oscillators-and-time:oscillator}{周期}给出的微小推力
使 $\theta_m$ 保持不变。
\begin{enumerate}[resume]
  \item 减半需要经过多少个\omterm{def:g12:oscillators-and-time:oscillator}{周期}？
  \item 在最大偏角处摆锤已升高 $h = L(1 - \cos\theta_m)
        \approx L\,\theta_m^2/2$（弧度）。推导 $E_m = \tfrac12
        m g L\,\theta_m^2$，并计算其值
        （见\cref{ch:g11:mechanical-energy}）。
  \item \omterm{def:g10:energy-conservation:energy}{能量}与\omterm{def:g12:oscillators-and-time:oscillator}{振幅}的平方成正比：每\omterm{def:g12:oscillators-and-time:oscillator}{周期}损失的\omterm{def:g10:energy-conservation:energy}{能量}占
        $E_m$ 的比例为多少？
  \item 推导每次推力提供的\omterm{def:g10:energy-conservation:energy}{能量}，以及所需的平均\omterm{def:g11:work-of-force:power}{功率}。
  \item 驱动重锤质量 $M = \qty{2.5}{kg}$，每周下降 $h =
        \qty{1.0}{m}$：这些\omterm{def:g10:energy-conservation:energy}{能量}是否足够？此时\omterm{def:g10:energy-conservation:efficiency}{效率}为多少？
\end{enumerate}

\medskip
\noindent\textbf{第四部分——石英挑战者。}一块石英机芯
以 $f = \qty{32768}{Hz}$ 振荡，每月约漂移
\qty{0.3}{s}；修复后的摆钟仍然每天
漂移 \qty{0.5}{s}。
\begin{enumerate}[resume]
  \item 求石英\omterm{def:g12:oscillators-and-time:oscillator}{周期}。证明 $\num{32768} = 2^{15}$，
        并说明为何 2 的幂次恰好符合手表电路的需要？
  \item 求石英表与修复后摆钟的相对漂移，并分别以每月秒数表示。
        谁更胜一筹？相差多少倍？
  \item 铯\omterm{def:g11:fundamental-interactions:ladder}{原子}钟标准——\num{9192631770} 个\omterm{def:g12:oscillators-and-time:oscillator}{周期}构成一秒——
        每一亿年漂移 \qty{1}{s}：其相对漂移为多少？
        比石英低多少个\omterm{def:g10:orders-of-magnitude:oom}{数量级}？
        为何 GPS 需要它？
  \item 结论以一句含数字的话给出：修复后的挂钟
        在墙上赢得了什么，现在由谁掌握“秒”的定义？
\end{enumerate}
\end{problem}
